
%%%%%%%%%%%%%%%%% 线性方程组 %%%%%%%%%%%%%%

\chapter{线性方程组}\label{chap:LinearSystems}

\section{线性方程组的多副面孔}

关于线性方程组为何物，有多种视角。第一种（也是最朴素的一种），就是将其视作含有$n$个未知数$x_1,x_2,\dots,x_n$的$m$个线性方程，
\[\left\{\begin{array}{cccccccccc}a_{1,1}x_1&+&a_{1,2}x_2&+&\ldots&+&a_{1,n}x_n&=&b_1\\a_{2,1}x_1&+&a_{2,2}x_2&+&\ldots&+&a_{2,n}x_n&=&b_2\\\cdots\\a_{m,1}x_1&+&a_{m,2}x_2&+&\ldots&+&a_{m,n}x_n&=&b_m&.\end{array}\right.\]

求解该方程组即求出{\bf 所有}同时满足这$m$个方程的数$x_1,x_2,\dots,x_n$构成的$n$-元组。

如果记$\mathbf{x}:=(x_{1},x_{2},\ldots,x_{n})^{T}\in\mathbb{F}^{n},\mathbf{b}=(b_{1},b_{2},\ldots,b_{m})^{T}\in\mathbb{F}^{m}$，且记
\[A=\begin{pmatrix}a_{1,1}&a_{1,2}&\ldots&a_{1,n}\\a_{2,1}&a_{2,2}&\ldots&a_{2,n}\\\vdots&\vdots&&\vdots\\a_{m,1}&a_{m,2}&\ldots&a_{m,n}\end{pmatrix},\]
那么上述线性方程组可以写成{\bf 矩阵形式}（也即{\bf 矩阵-向量方程}）
\[A\mathbf{x}=\mathbf{b}.\]
求解上述方程即求出所有满足$A\mathbf{x}=\mathbf{b}$的向量$\mathbf{x}\in\mathbb{F}^{n}$。

最后，回想一下矩阵-向量乘法的“列乘以坐标”规则，我们可以将原方程组写为{\bf 向量方程}
\[x_1\mathbf{a}_1+x_2\mathbf{a}_2+\cdots+x_n\mathbf{a}_n=\mathbf{b},\]
其中$\mathbf{a}_k$是矩阵$A$的第$k$列，$\mathbf{a}_k=(a_{1,k},a_{2,k},\dots,a_{m,k})^T,k=1,2,\dots,n$。

注意，上面三例本质上只是同种数学对象的不同表达而已。

在开始解释如何求解线性方程组之前，我们首先要注意，未知数的记法——比如$x_k,y_k$或者其他名称——都无关紧要。所以，求解线性方程组所必需的信息全部包含于两者：矩阵$A$（称为线性方程组的{\bf 系数矩阵}）和（右端）向量$\mathbf{b}$。进而，我们所需的全部信息都包含于以下矩阵
\pagebreak
\[\left(
  \begin{array}{cccc|c}
      a_{1,1} & a_{1,2} & \ldots & a_{1,n} & b_1    \\
      a_{2,1} & a_{2,2} & \ldots & a_{2,n} & b_2    \\
      \vdots  & \vdots  &        & \vdots  & \vdots \\
      a_{m,1} & a_{m,2} & \ldots & a_{m,n} & b_m    \\
    \end{array}
  \right)\]
该矩阵是通过将列向量$\mathbf{b}$增加到矩阵$A$右端获得的，称为线性方程组的{\bf 增广矩阵}。我们通常会在$A$和$\mathbf{b}$之间插入一条竖线来区分增广矩阵和系数矩阵。


\section{线性方程组的解、阶梯形和简化阶梯形}

线性方程组是用{\bf 高斯-若当消元法}（有时也称为{\bf 行化简}）来求解的。通过对线性方程组的增广矩阵的行作变换（也即对方程作变换），我们将其化为一简单的形式——即所谓{\bf 阶梯形}。当线性方程组具有{\bf 阶梯形}时，我们就很容易写出其解。

\subsection{行变换}

行变换有三种：
\begin{enumerate}
  \item 行交换：交换矩阵的两行；
  \item 缩放：将一行乘以非零标量$a$；
  \item 行置换：将第$k$行换为它与第$j$行的常数倍之和，而保持其他行不变。
\end{enumerate}

显然，变换1和2并不会改变线性方程组的解集——它们实质上就没有改变原方程组。

至于变换3，我们容易看出它不会使方程组的解减少。意即，令“新”方程组由“旧”方程组经过第3种行变换得来。那么“旧”方程组的每个解都是“新”方程组的解。

要看出为什么变换3不会导致解增加，也就是“新”方程组的任一解也是“旧”方程组的解，我们只需注意到第3种行变换是{\bf 可逆的}，也即，“旧”方程组也可以通过对“新”方程组进行第3种行变换得来（你能说出是什么样的变换吗？）

\subsubsection{行变换与乘以初等矩阵}

关于上述行变换为何不改变方程组的解集，还有一种更为“高级”的解释，也即，每种行变换都等价于左乘一种特殊的初等矩阵。

\setcounter{MaxMatrixCols}{12}
具体地说，乘以矩阵
\[
\begin{pNiceMatrix}[margin,first-row,first-col]
    &        &        &        & j      &        &        &        & k      &                        &              & \\
    & 1      &        &        & \vdots &        &        &        & \vdots & \Block{3-3}<\LARGE>{0} &              & \\
    &        & \ddots &        & \vdots &        &        &        & \vdots &                        & \hspace{1em} & \\
    &        &        & 1      & \vdots &        &        &        & \vdots &                        &              & \\
  j & \cdots & \cdots & \cdots & 0      & \cdots & \cdots & \cdots & 1      &                        &              & \\
    &        &        &        & \vdots & 1      &        &        & \vdots &                        &              & \\
    &        &        &        & \vdots &        & \ddots &        & \vdots &                        &              & \\
    &        &        &        & \vdots &        &        & 1      & \vdots &                        &              & \\
  k & \cdots & \cdots & \cdots & 1      & \cdots & \cdots & \cdots & 0      &                        &              & \\
    &        & \Block{3-2}<\LARGE>{0} &        &        &       &        &        &  &            1            &              & \\
    &        & \hspace{1em} &        &        &       &        &        &  &                        &        \ddots      & \\
    &        &  &        &        &       &        &        &  &                        &              & 1 
\end{pNiceMatrix}
\]
即交换第$j$行与第$k$行；乘以矩阵
\[
\begin{pNiceMatrix}[margin,first-col,first-row]
    &                        &        &   &   k    &   &                        &   \\
    & 1                      &        &   & \vdots &   & \Block{3-2}<\LARGE>{0} &   \\
    &                        & \ddots &   & \vdots &   &                        &   \\
    &                        &        & 1 & 0      &   &                        &   \\
  k & \cdots                 & \cdots & 0 & a      & 0 &                        &   \\
    & \Block{3-2}<\LARGE>{0} &        &   & 0      & 1 &                        &   \\
    &                        &        &   &        &   & \ddots                 & 0 \\
    &                        &        &   &        &   & 0                      & 1 \\
\end{pNiceMatrix}
\]
即将第$k$行乘以$a$。最后，乘以矩阵
\[
\begin{pNiceMatrix}[margin,first-row,first-col]
    &                        &        & j      &        & k      &                        &   \\
    & 1                      &        & \vdots &        & \vdots & \Block{2-2}<\LARGE>{0} &   \\
    &                        & \ddots & \vdots &        & \vdots &                        &   \\
  j & \cdots                 & \cdots & 1      & \cdots & 0      &                        &   \\
    &                        &        & \vdots & \ddots & \vdots &                        &   \\
  k & \cdots                 & \cdots & a      & \cdots & 1      &                        &   \\
    & \Block{2-2}<\LARGE>{0} &        &        &        &        & \ddots                 &   \\
    &                        &        &        &        &        &                        & 1
\end{pNiceMatrix}
\]
即将第$k$行加上第$j$行的$a$倍，并保持其他行不变。

要看出与这些矩阵相乘能实现上述效果，\marginpar{\small 一种描述（或记忆）这些初等矩阵的方式是：它们是向$I$施加对应行变换的结果。}只需要观察列向量乘以这些矩阵会发生什么变化。

注意所有初等矩阵都可逆（这对应于行变换的可逆性）。第一个矩阵的逆就是其本身。要得到第二个矩阵的逆，只需将$a$换成$1/a$。最后，通过将$a$换成$-a$可得第三个矩阵的逆。这些逆的正确性同样可通过观察列向量乘以这些矩阵会发生何种变化来检验。

因此，对方程组$A\mathbf{x}=\mathbf{b}$的增广矩阵作行变换等价于向该方程组左端乘上一个特殊的可逆矩阵$E$。在等式$A\mathbf{x}=\mathbf{b}$两端左乘$E$可得，方程\[A\mathbf{x}=\mathbf{b}\]的任意解也是方程\[EA\mathbf{x}=E\mathbf{b}\]的解。将这个方程左乘$E^{-1}$可得，其任意解是下述方程\[E^{-1}EA\mathbf{x}=E^{-1}E\mathbf{b}\]的解，也就是原方程$A\mathbf{x}=\mathbf{b}$的解。因此，行变换不改变方程组的解集。

\subsection{行化简}

行化简的主要步骤包含下面三步：
\begin{enumerate}
  \item 找到矩阵最左侧的非零列；
  \item 如有必要，通过施加第一类行变换（行交换），确保该列的第一个（最上方的）项是非零的。该项将被称为{\bf 主元}；
  \item 通过向第$2,3,\dots,m$行加（或减）第一行的适当倍数，“消去”主元之下的所有非零项（就是将它们都变为$0$）。
\end{enumerate}
我们先将这些步骤作用于整个矩阵，然后保持第一行不变并将这些步骤作用于第$2,3,\dots,m$行，然后是第$3,\dots,m$行，以此类推。

关键是记住在把某一行下方的所有行都减去该行的若干倍后（步骤3），我们就不再改变该行（甚至不会将该行与其他行进行交换）。

在有限次（至多$m$次）应用这些步骤后，我们就得到了矩阵的{\bf 阶梯形}。

\subsubsection{行化简实例}

考虑如下线性方程组：\[\left\{\begin{array}{c}x_1+2x_2+3x_3=1\\3x_1+2x_2+x_3=7\\2x_1+x_2+2x_3=1\end{array}\right.\]
该方程组的增广矩阵是
\[\left(
  \begin{array}{ccc|c}
      1 & 2 & 3 & 1   \\
      3 & 2 & 1 & 7   \\
      2 & 1 & 2 & 1   \\
    \end{array}
  \right)\]
将第二行减去第一行的$3$倍，并将第三行减去第一行的$2$倍，我们得到：
\[\begin{pNiceArray}{ccc|c}[margin,last-col]
  1 & 2 & 3 & 1 &  \\
  3 & 2 & 1 & 7 & -3R_1  \\
  2 & 1 & 2 & 1 & -2R_1 \\
\end{pNiceArray} \quad\sim\quad \left(
  \begin{array}{rrr|r}
      1 & 2 & 3 & 1   \\
      0 & -4 & -8 & 4   \\
      0 & -3 & -4 & -1   \\
    \end{array}
  \right)
\]
将第二行乘以$-1/4$可得
\[\left(
  \begin{array}{rrr|r}
    1 & 2 & 3 & 1   \\
    0 & 1 & 2 & -1   \\
    0 & -3 & -4 & -1   \\
  \end{array}
\right)\]
将第二行的$3$倍加到第三行，可得
\[\begin{pNiceArray}{rrr|r}[margin,last-col]
  1 & 2 & 3 & 1 &  \\
  0 & 1 & 2 & -1 &   \\
  0 & -3 & -4 & -1 & +3R_2 \\
\end{pNiceArray} \quad\sim\quad \left(
  \begin{array}{rrr|r}
      1 & 2 & 3 & 1   \\
      0 & 1 & 2 & -1   \\
      0 & 0 & 2 & -4   \\
    \end{array}
  \right)
\]
现在，我们就可用所谓{\bf 回代法}解此方程组了。也就是说，从最后一行（最后一个方程）得出$x_3 = -2$。接下来从第二个方程得到\[x_2=-1-2x_3=-1-2(-2)=3,\]
最后从第一行（第一个方程）得到\[x_1=1-2x_2-3x_3=1-6+6=1.\]
因此原方程组的解为\[\left\{\begin{array}{l}x_1=1,\\x_2=3,\\x_3=-2,\end{array}\right.\]
或用向量形式表示为\[\mathbf{x}=\begin{pmatrix}1\\3\\-2\end{pmatrix}\]
或写为$\mathbf{x}=(1,3,-2)^T$。我们可以通过作乘法$A\mathbf{x}$来检验该解，其中$A$为系数矩阵。

我们也可以不用回代法，而是从下到上做行变换，将系数矩阵主对角线上方的所有项全部变为$0$：首先将最后一行乘以$1/2$，后续过程则不言而喻：
\begin{align*}
  \begin{pNiceArray}{rrr|r}[margin,last-col]
    1 & 2 & 3 & 1 & -3R_3 \\
    0 & 1 & 2 & -1 &  -2R_3 \\
    0 & 0 & 1 & -2 &  \\
  \end{pNiceArray} &\sim 
  \begin{pNiceArray}{rrr|r}[margin,last-col]
        1 & 2 & 0 & 7 & -2R_2 \\
        0 & 1 & 0 & 3   \\
        0 & 0 & 1 & -2   \\
  \end{pNiceArray}\\
  &\sim\begin{pNiceArray}{rrr|r}[margin]
    1 & 0 & 0 & 1 \\
    0 & 1 & 0 & 3   \\
    0 & 0 & 1 & -2   \\
\end{pNiceArray}
\end{align*}
然后从增广矩阵一眼可知解为$\mathbf{x}=(1,3,-2)^T$。

留给读者一道习题：梳理从下到上阶段的行化简算法。

\subsection{阶梯形}\label{subsec:echelon}

矩阵具有{\bf 阶梯形}，若其满足如下两个条件：
\begin{enumerate}
  \item 所有的零行（即全部的项都等于$0$的行）位于所有非零项的下方；
\end{enumerate}
对于非零行，我们称最左边的非零项为{\bf 先导元}。那么阶梯形的第二个性质可表述为：
\begin{enumerate}[start=2]
  \item 任意非零行的先导元必定严格处于前一行先导元的右侧。
\end{enumerate}

阶梯形中每一行的先导元也称为{\bf 主元}，\marginpar{\small 主元：一行中的先导元（最左侧的非零项）。}因为这些项恰为行化简时所用的{\bf 主元}。

阶梯形有个特例是所谓{\bf 三角}形。在上面的例子中我们就得到了这种形式。在这种形式中，系数矩阵是方阵（$n\times n$矩阵），其主对角线项均不为零，而主对角线之下的所有项都为零。右端项，或者说增广矩阵的最右列可以是任意的。

在行化简的从下到上阶段之后，我们就得到了矩阵的所谓{\bf 简化阶梯形}：系数矩阵等于$I$（如上例所示）是简化阶梯形的一个特例。

一般的定义是：称矩阵具有{\bf 简化阶梯形}，若其具有阶梯形且
\begin{enumerate}[start=3]
  \item 所有主元都等于$1$；
  \item 主元之上的所有项都为$0$。注意，主元之下的所有项也都是$0$，因为该矩阵具有阶梯形。
\end{enumerate}

为了从阶梯形得到简化阶梯形，我们按从下到上、从右到左的次序，利用行置换消去主元之上的所有项。

简化阶梯形的一个例子是系数矩阵等于$I$的方程组。在这种情况下，从简化阶梯形就很容易看出解。在一般的情形下，从简化阶梯形看出解也很容易。例如，令方程组（增广矩阵）的简化阶梯形是\[\left(\begin{array}{ccccc|c}\boxed{1}&2&0&0&0&1\\0&0&\boxed{1}&5&0&2\\0&0&0&0&\boxed{1}&3\end{array}\right);\]
此处我们将主元框出来了。求解思路是，将没有主元的列所对应的变量（所谓{\bf 自由变量}）移到右侧。然后我们可直接写出解为\pagebreak\[\left\{\begin{array}{l}x_1=1-2x_2\\x_2~\text{为自由变量}\\x_3=2-5x_4\\x_4~\text{为自由变量}\\x_5=3\end{array}\right.\]
或者写成向量形式
\[\mathbf{x}=\begin{pmatrix}1-2x_2\\x_2\\2-5x_4\\x_4\\3\end{pmatrix}=\begin{pmatrix}1\\0\\2\\0\\3\end{pmatrix}+x_2\begin{pmatrix}-2\\1\\0\\0\\0\end{pmatrix}+x_4\begin{pmatrix}0\\0\\-5\\1\\0\end{pmatrix},\quad x_2,x_4\in\mathbb{F}.\]

也可用回代法从阶梯形求出解：基本思路是从下到上操作，将所有自由变量移到右侧。
\begin{problemset}
  \item 将下列线性方程组写成矩阵形式和向量方程形式：
  \begin{enumerate}[label={\alph*)}]
    \item $\left\{\begin{array}{rrrrrrrr}x_1&+&2x_2&-&x_3&=&-1\\2x_1&+&2x_2&+&x_3&=&1\\3x_1&+&5x_2&-&2x_3&=&-1\end{array}\right.$
    \item $\left\{\begin{array}{rrrrrrrr}x_1&-&2x_2&-&x_3&=&1\\2x_1&-&3x_2&+&x_3&=&6\\3x_1&-&5x_2&&&=&7\\x_1&&&+&5x_3&=&9\end{array}\right.$
    \item $\left\{\begin{array}{rrrrrrrrrr}x_1&+&2x_2&&&+&2x_4&=&6\\3x_1&+&5x_2&-&x_3&+&6x_4&=&17\\2x_1&+&4x_2&+&x_3&+&2x_4&=&12\\2x_1&&&-&7x_3&+&11x_4&=&7\end{array}\right.$
    \item $\left\{\begin{array}{rrrrrrrrrr}x_1&-&4x_2&-&x_3&+&x_4&=&3\\2x_1&-&8x_2&+&x_3&-&4x_4&=&9\\-x_1&+&4x_2&-&2x_3&+&5x_4&=&-6\end{array}\right.$
    \item $\left\{\begin{array}{rrrrrrrrrr}x_1&+&2x_2&-&x_3&+&3x_4&=&2\\2x_1&+&4x_2&-&x_3&+&6x_4&=&5\\&&x_2&&&+&2x_4&=&3\end{array}\right.$
    \item $\left\{\begin{array}{rrrrrrrrrrr}2x_1&-&2x_2&-&x_3&+&6x_4&-&2x_5&=&1\\x_1&-&x_2&+&x_3&+&2x_4&-&x_5&=&2\\4x_1&-&4x_2&+&5x_3&+&7x_4&-&x_5&=&6\end{array}\right.$
    \item $\left\{\begin{array}{rrrrrrrrrrrr}3x_1&-&x_2&+&x_3&-&x_4&+&2x_5&=&5\\x_1&-&x_2&-&x_3&-&2x_4&-&x_5&=&2\\5x_1&-&2x_2&+&x_3&-&3x_4&+&3x_5&=&10\\2x_1&-&x_2&&&-&2x_4&+&x_5&=&5\end{array}\right.$
  \end{enumerate}
  解这些方程组。将解写成向量形式。
  \item 求出向量方程\[x_1\mathbf{v}_1+x_2\mathbf{v}_2+x_3\mathbf{v}_3=\mathbf{0},\]的{\bf 全部}解，其中$\mathbf{v}_{1}=(1,1,0)^{T},\mathbf{v}_{2}=(0,1,1)^{T},\mathbf{v}_{3}=(1,0,1)^{T}$。关于向量组$\mathbf{v}_1,\mathbf{v}_2,\mathbf{v}_3$的线性无关性（线性相关性），你可以得出什么结论？
\end{problemset}

\section{主元分析}

通过分析线性方程组增广矩阵的阶梯形（简化阶梯形）的主元，可以回答所有关于解的存在性和唯一性的问题。首先要分析的问题是：方程$A\mathbf{x}=\mathbf{b}$何时会{\bf 不相容}（无解）。稍加思索可立得答案：
\begin{center}
  \fbox{\parbox{0.8\linewidth}{方程组不相容（无解），当且仅当{\bf 增广}矩阵的阶梯形的最后一列存在主元，也就是，当且仅当增广矩阵的阶梯形有一行是$\left(\begin{array}{cccc|c}0&0&\ldots&0&b\end{array}\right),b\neq0$。}}
\end{center}
的确如此：这样一行对应于方程$0x_1+0x_2+\ldots+0x_n=b\neq0$，而该方程无解。如果不存在这样的行，我们就可以将增广矩阵化为简化阶梯形并直接看出方程组的解。

下面再列出三个结论。注意，它们是与方程组的{\bf 系数矩阵}而非增广矩阵相关。
\begin{enumerate}
  \item 方程组的解（若存在）是唯一的，当且仅当没有自由变量，也就是当且仅当系数矩阵的阶梯形在每一列中都有主元；
  \item 方程$A\mathbf{x}=\mathbf{b}$对于任意右端项$\mathbf{b}$都是相容的，当且仅当系数矩阵的阶梯形在每一行中都有主元；
  \item 方程$A\mathbf{x}=\mathbf{b}$对于任意右端项$\mathbf{b}$都有{\bf 唯一解}，当且仅当系数矩阵$A$的阶梯形在每一行和每一列中都有主元。
\end{enumerate}

第一个结论很平凡，因为自由变量对应于非唯一性。只是要强调，这一结论{\bf 并不保证}存在性。

第二个结论稍复杂些。若系数矩阵的每行都有主元，{\bf 增广}矩阵的最后一列就不会有主元，因此无论右端项$\mathbf{b}$是什么，方程组都是相容的。

下面证明，如果系数矩阵$A$的阶梯形中有零行，那么可以选择一个右端项$\mathbf{b}$使得方程组$A\mathbf{x}=\mathbf{b}$是不相容的。令$A_{\mathrm{e}}$为系数矩阵$A$的阶梯形。那么\[A_{\mathrm{e}}=EA\]其中$E$是一系列初等矩阵之积，对应于一系列行变换，$E=E_N\dots E_2E_1$。若$A_{\mathrm{e}}$有零行，那么其最后一行也为零。因此，如果取$\mathbf{b}_\mathrm{e}=(0,\ldots,0,1)^T$（除了最后一项以外各项全为$0$）那么方程\[A_{\mathrm{e}}\mathbf{x}=\mathbf{b}_e\]无解。方程两端同时左乘$E^{-1}$，并运用$E^{-1}A_{\mathrm{e}}=A$，可得方程\[A\mathbf{x}=E^{-1}\mathbf{b}_\mathrm{e}\]
无解。

最后，由结论1和2可立得结论3。{\hfill\qedsymbol}

从以上对于主元的分析可得若干重要推论。为此我们需要注意到
\begin{center}
  \fbox{\parbox{0.8\linewidth}{在阶梯形中，任一行和任一列中至多存在一个主元（可以没有主元）。}}
\end{center}

\subsection{有关线性无关性和基的推论、维数}

通过行化简，很容易回答$\mathbb{F}^n$中的向量组何时为基、线性无关组或张成组的问题。

\begin{proposition}\label{prop:li_span_eq}
  设有向量组$\mathbf{v}_{1},\mathbf{v}_{2},\ldots,\mathbf{v}_{m}\in\mathbb{F}^{n}$，并令$A=[\mathbf{v}_{1},\mathbf{v}_{2},\ldots,\mathbf{v}_{m}]$是各列为$\mathbf{v}_{1},\mathbf{v}_{2},\ldots,\mathbf{v}_{m}$的$n\times m$矩阵。那么
  \begin{enumerate}
    \item 向量组$\mathbf{v}_{1},\mathbf{v}_{2},\ldots,\mathbf{v}_{m}$是线性无关的，当且仅当$A$的阶梯形的每一列都有主元；
    \item 向量组$\mathbf{v}_{1},\mathbf{v}_{2},\ldots,\mathbf{v}_{m}$是$\mathbb{F}^{n}$中的完备组（张成组、生成组），当且仅当$A$的阶梯形的每一行都有主元；
    \item 向量组$\mathbf{v}_{1},\mathbf{v}_{2},\ldots,\mathbf{v}_{m}$是$\mathbb{F}^{n}$中的基，当且仅当$A$的阶梯形的每一行与每一列都有主元。
  \end{enumerate}
\end{proposition}
\begin{proof}
  向量组$\mathbf{v}_{1},\mathbf{v}_{2},\ldots,\mathbf{v}_{m}\in\mathbb{F}^{n}$线性无关，当且仅当方程\[x_1\mathbf{v}_1+x_2\mathbf{v}_2+\ldots+x_m\mathbf{v}_m=\mathbf{0}\]具有唯一（平凡）解$x_{1}=x_{2}=\ldots=x_{m}=0$，或等价于，方程$A\mathbf{x}=\mathbf{0}$具有唯一解$\mathbf{x}=\mathbf{0}$。根据前述结论1，这些等价于矩阵的阶梯形每一列都有主元。

  类似地，向量组$\mathbf{v}_{1},\mathbf{v}_{2},\ldots,\mathbf{v}_{m}\in\mathbb{F}^{n}$是$\mathbb{F}^{n}$中的完备组，当且仅当方程\[x_1\mathbf{v}_1+x_2\mathbf{v}_2+\cdots+x_m\mathbf{v}_m=\mathbf{b}\]对任意右端项$\mathbf{b}\in\mathbb{F}^n$都有解。根据前述结论2，这等价于矩阵的阶梯形的每一行都有主元。

  最后，向量组$\mathbf{v}_{1},\mathbf{v}_{2},\ldots,\mathbf{v}_{m}\in\mathbb{F}^{n}$是$\mathbb{F}^{n}$中的基，当且仅当方程\[x_1\mathbf{v}_1+x_2\mathbf{v}_2+\cdots+x_m\mathbf{v}_m=\mathbf{b}\]对任意右端项$\mathbf{b}\in\mathbb{F}^n$都有唯一解。根据前述结论3，这等价于$A$的阶梯形的每一行和每一列都有主元。
\end{proof}

\begin{proposition}\label{prop:li_no_more}
  $\mathbb{F}^{n}$中的任意线性无关向量组不能含有超过$n$个向量。
\end{proposition}

\begin{proof}
  设向量组$\mathbf{v}_{1},\mathbf{v}_{2},\ldots,\mathbf{v}_{m}\in\mathbb{F}^{n}$线性无关，并令$A=[\mathbf{v}_{1},\mathbf{v}_{2},\ldots,\mathbf{v}_{m}]$是各列为$\mathbf{v}_{1},\mathbf{v}_{2},\ldots,\mathbf{v}_{m}$的$n\times m$矩阵。由命题 \ref{prop:li_span_eq}，$A$的阶梯形的每列都必有主元，而这在$m>n$时是不可能的（主元的数目不会超过行数）。
\end{proof}

\begin{proposition}\label{prop:basis_same_size}
  向量空间$V$中的任意两个基都有相同个数的向量。
\end{proposition}

\begin{proof}
  设$\mathbf{v}_{1},\mathbf{v}_{2},\ldots,\mathbf{v}_{n}$和$\mathbf{w}_{1},\mathbf{w}_{2},\ldots,\mathbf{w}_{m}$是$V$中两个不同的基。不失一般性，可假设$n\le m$。考虑定义如下的同构$A:\mathbb{F}^{n}\to V$：\[A\mathbf{e}_k=\mathbf{v}_k,\quad k=1,2,\ldots n,\]其中$\mathbf{e}_1,\mathbf{e}_2,\ldots,\mathbf{e}_n$是$\mathbb{R}^n$中的标准基。

  由于$A^{-1}$也是同构，因此向量组\[A^{-1}\mathbf{w}_1,A^{-1}\mathbf{w}_2,\ldots,A^{-1}\mathbf{w}_m\]也是基（见第 \ref{chap:BasicNotion} 章的\autoref{thm:preserve_basis}）。因此它是线性无关的，则由命题 \ref{prop:li_no_more} 可知$m\le n$。结合$n\le m$的假设即得$m=n$。
\end{proof}

接下来的结论是上述命题的一个特殊情况。
\begin{proposition}\label{prop:basis_size_n}
    $\mathbb{F}^{n}$中的任意一个基都恰有$n$个向量。
\end{proposition}
\begin{proof}
  从上个命题可立即得到此结论，不过也可以直接证明。设$\mathbf{v}_{1},$ $\mathbf{v}_{2},\ldots,\mathbf{v}_{m}$是$\mathbb{F}^{n}$中的基，并令$A$是各列为$\mathbf{v}_{1},\mathbf{v}_{2},\ldots,\mathbf{v}_{m}$的$n\times m$矩阵。该向量组是基，意味着方程\[A\mathbf{x}=\mathbf{b}\]对任意（所有可能的）右端项$\mathbf{b}$都有唯一解。解的存在性表明该矩阵的简化阶梯形每行都有主元，因此主元的数目恰为$n$，唯一性则表明系数矩阵（的阶梯形）每列都有主元。因此\[m = \text{列数} = \text{主元数} = n\]
\end{proof}

\begin{proposition}\label{prop:span_more}
  $\mathbb{F}^{n}$中的任意张成（生成）组都至少包含$n$个向量。
\end{proposition}
\begin{proof}
  设$\mathbf{v}_{1},$ $\mathbf{v}_{2},\ldots,\mathbf{v}_{m}$是$\mathbb{F}^{n}$中的完备组，并令$A$是各列为$\mathbf{v}_{1},\mathbf{v}_{2},\ldots,$ $\mathbf{v}_{m}$的$n\times m$矩阵。命题 \ref{prop:li_span_eq} 的结论2表明，$A$的阶梯形的每行都有主元。因为主元数目不超过列数，所以$n\le m$。
\end{proof}

\subsection{有关可逆矩阵的推论}

\begin{proposition}
  矩阵$A$可逆，当且仅当其阶梯形的每一行和每一列都有主元。
\end{proposition}

\begin{proof}
  我们在本节开头已讨论过，方程$A\mathbf{x}=\mathbf{b}$对于任意右端项$\mathbf{b}$都有唯一解，当且仅当$A$的阶梯形在每一行和每一列中都有主元。又由第 \ref{chap:BasicNotion} 章定理 \ref{thm:invert_unique} 知，矩阵（线性变换）$A$是可逆的，当且仅当方程$A\mathbf{x}=\mathbf{b}$对于任意右端项$\mathbf{b}$都有唯一解。

  还有另一种证法：我们知道矩阵可逆当且仅当它的各列形成一个基（见第 \ref{chap:BasicNotion} 章 \ref{subsec:invert_eq} 小节的推论 \ref{col:invert_basis}）。而上述命题 \ref{prop:basis_size_n} 表明这等价于矩阵的阶梯形每一行和每一列都有主元。
\end{proof}

由上述命题立得下述结论：
\begin{corollary}
  可逆矩阵必为方阵（$n\times n$矩阵）。
\end{corollary}
\begin{proposition}
  若方阵（$n\times n$矩阵）是左可逆的或右可逆的，那么它是可逆的。换言之，为检验方阵$A$的可逆性，只需检验$AA^{-1}=I$, $A^{-1}A=I$之一。
\end{proposition}

注意，此命题只适用于方阵！

\begin{proof}
  我们知道矩阵$A$可逆当且仅当方程$A\mathbf{x}=\mathbf{b}$对于任意右端项$\mathbf{b}$都有唯一解。而这又等价于矩阵$A$的阶梯形在每一行和每一列中都有主元。

  若矩阵$A$是左可逆的，那么方程$A\mathbf{x}=\mathbf{0}$有唯一解$\mathbf{x}=\mathbf{0}$。的确如此：若$B$是$A$的左逆（也即$BA=I$），而$\mathbf{x}$满足\[A\mathbf{x}=\mathbf{0},\]
  那么将此等式两端同时左乘$B$，即得到$\mathbf{x}=\mathbf{0}$，则此解唯一。因此，$A$的阶梯形在每列中都有主元（无自由变量）。若矩阵$A$是方阵（$n\times n$矩阵），那么其阶梯形的每行也必有主元（共$n$个主元，而且每行至多有$1$个主元），因此该矩阵是可逆的。

  若矩阵$A$是右可逆的，且$C$是$A$的右逆（$AC=I$），那么对于$\mathbf{x}=C\mathbf{b},\mathbf{b}\in\mathbb{F}^{n}$，有\[A\mathbf{x}=AC\mathbf{b}=I\mathbf{b}=\mathbf{b}.\]
  因此，对任意右端项$\mathbf{b}$，方程$A\mathbf{x}=\mathbf{b}$都有解$\mathbf{x}=C\mathbf{b}$。于是$A$的阶梯形每一行都有主元。若$A$是方阵，那么它的每一列也都有主元，故$A$可逆。
\end{proof}

\begin{problemset}
  \item $b$取何值时，方程组\[\begin{pmatrix}1&2&2\\2&4&6\\1&2&3\end{pmatrix}\mathbf{x}=\begin{pmatrix}1\\4\\b\end{pmatrix}\]有解？对于该$b$，求方程组的通解。
  \item 判断向量组\[\begin{pmatrix}1\\1\\0\\0\end{pmatrix},\quad\begin{pmatrix}1\\0\\1\\0\end{pmatrix},\quad\begin{pmatrix}0\\0\\1\\1\end{pmatrix},\quad\begin{pmatrix}0\\1\\0\\1\end{pmatrix}\]是否为线性无关的。
  
  这四个向量能张成$\RR^4$吗？（换言之，它是生成组吗？）$\CC^4$呢？
  \item 判断下面哪些向量组是$\RR^3$中的基？
  \begin{enumerate}[label={\alph*)}]
    \item $(1,2,-1)^{T},(1,0,2)^{T},(2,1,1)^{T}$；
    \item $(-1,3,2)^{T},(-3,1,3)^{T},(2,10,2)^{T}$；
    \item $(67,13,-47)^{T},(\pi,-7.84,0)^{T},(3,0,0)^{T}$。
  \end{enumerate}
  哪些向量组是$\CC^3$的基？
  \item 多项式$x^{3}+2x,x^{2}+x+1,x^{3}+5$能否生成（张成）$\mathbb{P}_3$？证明你的结论。
  \item $\mathbb{F}^4$中的$5$个向量是否线性无关？证明你的结论。
  \item 证明或证伪：若方阵（$n\times n$矩阵）$A$的列是线性无关的，那么$A^2=AA$的列也是。
  \item 证明或证伪：若方阵（$n\times n$矩阵）$A$的列是线性无关的，那么$A^3=AAA$的列也是。
  \item 证明：如果方程$A\mathbf{x}=\mathbf{0}$有唯一解（即$A$的阶梯形每一列都有主元），那么$A$是左可逆的。
  
  {\bf 提示}：初等矩阵可能有用处。
  
  {\bf 注意}：正文中已经证明，若$A$是左可逆的则方程$A\mathbf{x}=\mathbf{0}$具有唯一解。此处你需要证明的是该结论的逆命题，而正文中并未给出证明。

  {\bf 注}：这题可能很难，因为它需要对所论述的主题有深刻的理解。然而，当理解需要做些什么之后，这道题就非常简单了。
  \pagebreak
  \item 矩阵的简化阶梯形是唯一的吗？证明你的结论。
  
  即，假设我们通过作任意行变换（不按照特定的算法）得到简化阶梯形矩阵。究竟是每次都会得到相同的矩阵，还是可能得到不同的矩阵？注意，只能作行变换，不允许做列变换。

  {\bf 提示}：如果你是从可逆矩阵开始做变换呢？另外，主元是始终在同样的列中，还是根据所做的行变换不同而处于不同的列？如果你可以不通过（对简化阶梯形）作相反的行变换来推知主元列的位置，那么主元列的位置就不取决于行变换。
\end{problemset}

\section{通过行化简求\texorpdfstring{$A^{-1}$}{A**\{-1\}}}

如上面所讨论的，可逆矩阵必为方阵，并且其阶梯形必定在每一行和每一列中都有主元。所以可逆矩阵的简化阶梯形是恒等矩阵$I$。因此
\begin{center}
  \fbox{\parbox{0.8\linewidth}{任意可逆矩阵都行等价于（即可通过行变换化为）恒等矩阵。}}
\end{center}

下面陈述一个用于求出$n\times n$矩阵的逆的简便算法：
\begin{enumerate}
  \item 向$A$的右侧附上$n\times n$恒等矩阵，将其{\bf 扩充}成$n\times 2n$的矩阵$(A|I)$；
  \item 对扩充后的矩阵作行变换，将$A$变换为恒等矩阵$I$；
  \item 我们添加的矩阵$I$在此过程中自动转化为$A^{-1}$；
  \item 如果通过行变换不能把$A$变换为恒等矩阵，那么$A$不可逆。
\end{enumerate}

上述算法有好几种解释。第一种（也是比较朴素的一种）是：我们知道对于可逆矩阵$A$，向量$A^{-1}\mathbf{b}$是方程$A\mathbf{x}=\mathbf{b}$的解。所以为了求出$A^{-1}$的第$k$列，我们需要求出$A\mathbf{x}=\mathbf{e}_k$的解，其中$\mathbf{e}_1,\mathbf{e}_2,\ldots,\mathbf{e}_n$是$\mathbb{R}^n$中的标准基。上述算法就是同时求解了如下$n$个方程\[A\mathbf{x}=\mathbf{e}_k,\quad k=1,2,\ldots,n\]

我们再展示一种更为“高级”的解释。之前我们讨论过，每次行变换都可通过左乘初等矩阵来实现。令$E_{1},E_{2},\ldots,E_{N}$是对应于我们所作的行变换的初等矩阵，并令$E=E_{N}\cdots E_{2}E_{1}$为它们的乘积\footnote{尽管此处顺序无关紧要，但请注意，若行变换$E_1$是最先进行的，那么在乘积中$E_1$是最右侧的一项。}。我们知道行变换可将$A$变为恒等矩阵，也即$EA=I$，故$E=A^{-1}$。同样的行变换作用于扩充的矩阵$(A|I)$，会将其变为$(EA|E) = (I|A^{-1})$。{\hfill \qedsymbol}

这种运用初等矩阵的“高级”解释蕴涵了之后会用到的一个重要命题：

\begin{theorem}\label{thm:invert_elementary_prod}
  任意可逆矩阵都可表示成若干初等矩阵之积。
\end{theorem}

\begin{proof}
  如前两段所讨论的，$A^{-1}=E_{N}\cdots E_{2}E_{1}$，所以\[A=(A^{-1})^{-1}=E_{1}^{-1}E_{2}^{-1}\cdots E_{N}^{-1}.\]
\end{proof}

\begin{example*}
  假设我们要求出以下矩阵的逆。
  \[\left(\begin{array}{rrr}1&4&-2\\-2&-7&7\\3&11&-6\end{array}\right).\]
  向其右侧附上恒等矩阵，并使用行变换可得
  \begin{align*}
  &\begin{pNiceArray}{rrr|rrr}[last-col]
    1&4&-2&1&0&0&\\
    -2&-7&7&0&1&0&+2R_1\\
    3&11&-6&0&0&1&-3R_1
  \end{pNiceArray}&&\sim 
  \begin{pNiceArray}{rrr|rrr}[last-col]
    1&4&-2&1&0&0&\\
    0&1&3&2&1&0&\\
    0&-1&0&-3&0&1&+R_2
  \end{pNiceArray}\sim\\
  &\begin{pNiceArray}{rrr|rrr}[last-col]
    1&4&-2&1&0&0&\times 3\\
    0&1&3&2&1&0&\\
    0&0&3&-1&1&1&
\end{pNiceArray}&&\sim 
  \begin{pNiceArray}{rrr|rrr}[last-col]
    3&12&-6&3&0&0&+2R_3\\
    0&1&3&2&1&0&-R_3\\
    0&0&3&-1&1&1&
  \end{pNiceArray}\sim
\end{align*}
在最后一次行变换中，我们将第一行乘以$3$倍是为了在行化简的从下到上阶段中不出现分数。
继续作行变换有
\begin{align*}
  \begin{pNiceArray}{rrr|rrr}[margin,last-col]
    3&12&0&1&2&2&-12R_2\\
    0&1&0&3&0&-1&\\
    0&0&3&-1&1&1&
  \end{pNiceArray}&\sim 
  \begin{pNiceArray}{rrr|rrr}[margin]
    3&0&0&-35&2&14\\
    0&1&0&3&0&-1\\
    0&0&3&-1&1&1
  \end{pNiceArray}
\end{align*}
将第一行和第三行除以$3$，就得到了逆矩阵\[\begin{pmatrix}-35/3&2/3&14/3\\3&0&-1\\-1/3&1/3&1/3\end{pmatrix}\]
\end{example*}

\begin{problemset}
  \item 求出下列矩阵的逆：
  \[\left(\begin{array}{ccc}1&2&1\\3&7&3\\2&3&4\end{array}\right),\quad\left(\begin{array}{ccc}1&-1&2\\1&1&-2\\1&1&4\end{array}\right).\]
  给出所有的步骤。
\end{problemset}

\section{维数、有限维向量空间}

\begin{definition*}
  向量空间$V$的维数$\dim V$是基所具有的向量个数。

  对于只包含零向量$\mathbf{0}$的向量空间，规定$\dim V=0$；若$V$不存在（有限的）基，则规定$\dim V=\infty$。

  若$\dim V$是有限值，则称$V$是{\bf 有限维的}；否则称其为{\bf 无限维的}。
\end{definition*}

命题 \ref{prop:basis_same_size} 确保了维数是定义完善的，也即，它并不依赖于基的选取。

第 \ref{chap:BasicNotion} 章的命题 \ref{prop:generating_contain_basis} 表明，向量空间$V$的任意有限张成组都包含一个基。由此立得
\begin{proposition}
  向量空间$V$是有限维的，当且仅当其包含有限张成组。
\end{proposition}

假设我们想检验有限维向量空间中的某向量组是否为基（或者检验其是否为线性无关的，或检验其是否完备）。最简单的方法大概是，利用同构$A:V\to \mathbb{R}^n,n=\dim V$将原问题转化到$\mathbb{R}^n$中，然后就可用行化简（分析主元）来回答这样的问题。

注意，若$\dim V=n$，那么同构$A:V\to \mathbb{R}^n$总是存在的。的确如此：若$\dim V=n$，那么存在基$\mathbf{v}_{1},\mathbf{v}_{2},\ldots,\mathbf{v}_{n}\in V$，于是可定义同构$A:V\to \mathbb{R}^n$为
\[A\mathbf{v}_k=\mathbf{e}_k,\quad k=1,2,\ldots,n.\]

下面给出命题 \ref{prop:li_no_more} 和 \ref{prop:span_more} 的两个推论，作为上述方法的示例。
\begin{proposition}\label{prop:arbit_li_no_more}
  有限维向量空间$V$中的任意线性无关组不会具有超过$\dim V$个向量。
\end{proposition}
\begin{proof}
  令$\mathbf{v}_1,\mathbf{v}_2,\ldots,\mathbf{v}_m\in V$为线性无关组，并令$A:V\to \mathbb{R}^n$为同构。那么$A\mathbf{v}_{1},A\mathbf{v}_{2},\ldots,A\mathbf{v}_{m}$是$\mathbb{R}^n$中的线性无关组，则由命题 \ref{prop:li_no_more} 知$m\le n$。
\end{proof}
\begin{proposition}\label{prop:arbit_span_more}
  有限维向量空间$V$中的任意生成组一定具有至少$\dim V$个向量。
\end{proposition}
\begin{proof}
  令$\mathbf{v}_1,\mathbf{v}_2,\ldots,\mathbf{v}_m\in V$为完备组，并令$A:V\to \mathbb{R}^n$为同构。那么$A\mathbf{v}_{1},A\mathbf{v}_{2},\ldots,A\mathbf{v}_{m}$是$\mathbb{R}^n$中的完备组，则由命题 \ref{prop:span_more} 知$m\ge n$。
\end{proof}

\subsection{将线性无关组补足为基}

下面结论将在后文发挥重要作用。

\begin{proposition}[补足为基]
  有限维向量空间中的线性无关组可被补足为基，也即，若$\mathbf{v}_1,\mathbf{v}_2,\ldots,\mathbf{v}_r$是有限维向量空间$V$中的线性无关向量组，那么可以找到向量$\mathbf{v}_{r+1},\mathbf{v}_{r+2}\ldots,\mathbf{v}_{n}$，使得向量组$\mathbf{v}_1,\mathbf{v}_2,\ldots,\mathbf{v}_n$是$V$中的基。
\end{proposition}

\begin{proof}
  令$n=\dim V$。取任意不属于$\operatorname{span}\{\mathbf{v}_{1},\mathbf{v}_{2},\ldots,\mathbf{v}_{r}\}$的向量，记其为$\mathbf{v}_{r+1}$（这总是可以完成的，因为向量组$\mathbf{v}_1,\mathbf{v}_2,\ldots,\mathbf{v}_r$不是生成组）。由第 \ref{chap:BasicNotion} 章的习题 \ref{ex:1_2_5}，向量组$\mathbf{v}_1,\mathbf{v}_2,\ldots,\mathbf{v}_r,\mathbf{v}_{r+1}$是线性无关的（注意，由命题 \ref{prop:arbit_li_no_more}，$r<n$）。对新构造的向量组重复此步骤，可得到向量$\mathbf{v}_{r+2}$，以此类推。

  当我们得到生成组时，就停止此过程。注意，这一过程不会无限持续下去，因为$V$中的线性无关组不会具有超过$n=\dim V$个向量。
\end{proof}

\subsection{有限维向量空间的子空间}

\begin{theorem}\label{thm:subspace_finite_dimensional}
  令$V$是向量空间$W$的子空间，且$\dim W < \infty$。那么$V$是有限维的，且$\dim V\le \dim W$。

  此外，若$\dim V= \dim W$，则$V=W$（仍假设$V$是$W$的子空间）。
\end{theorem}
\begin{remark*}
  这定理初看上去没什么意思——似乎只是命题 \ref{prop:arbit_li_no_more} 的简单推论。然而，仅在已有$V$的基时，才能应用命题 \ref{prop:arbit_li_no_more}。但此处我们只有$W$中的基，并且无法确定该基中有多少向量属于$V$。事实上，$W$的这个基中可能没有一个向量属于$V$——很容易举出这样的例子。
\end{remark*}

\begin{proof}[\autoref{thm:subspace_finite_dimensional} 的证明]
  $V = \{\mathbf{0}\}$为该定理的平凡情形。现在假设$V \ne \{\mathbf{0}\}$。

  我们的目标是找到$V$的一个基。取非零向量$\mathbf{v}_1\in V$。若$V=\operatorname{span}\{\mathbf{v}_1\}$，我们就得到了所需的基（包含单一向量$\mathbf{v}_1$）。

  如若不然，则用归纳法。假设我们已经构造了$r$个线性无关的向量$\mathbf{v}_{1},\ldots,\mathbf{v}_{r}\in V$。若$V=\mathrm{span}\{\mathbf{v}_{k}:1\leq k\leq r\}$，那么我们就寻得了$V$的基。否则，就存在向量$\mathbf{v}_{r+1}\in V$，满足$\mathbf{v}_{r+1}\notin\mathrm{span}\{\mathbf{v}_{k}:1\leq k\leq r\}$。由第 \ref{chap:BasicNotion} 章的习题 \ref{ex:1_2_5}，向量组$\mathbf{v}_1,\mathbf{v}_2,\ldots,\mathbf{v}_r,\mathbf{v}_{r+1}$是线性无关的。
\end{proof}

\begin{problemset}
  \item 判断正误：
  \begin{enumerate}[label={\alph*)}]
    \item 每个由有限向量组生成的向量空间都有基；
    \item 每个向量空间都有（有限）基；
    \item 一向量空间不会含有超过一个基；
    \item 若一向量空间包含有限基，那么每个基所具有的向量数目都相同；
    \item $\mathbb{P}_{n}$的维数是$n$；
    \item $M_{m\times n}$的维数是$m+n$；
    \item 若向量$\mathbf{v}_{1},\mathbf{v}_{2},\ldots,\mathbf{v}_{n}$生成（张成）向量空间$V$，那么$V$中每个向量都可唯一写成向量$\mathbf{v}_{1},\mathbf{v}_{2},\ldots,\mathbf{v}_{n}$的线性组合；
    \item 有限维向量空间的每个子空间都是有限维的；
    \item 如果$V$是维数为$n$的向量空间，那么$V$恰有一个维数为$0$和一个维数为$n$的子空间。
  \end{enumerate}
  \item 证明：若$V$是维数为$n$的向量空间，那么$V$中的向量组$\mathbf{v}_{1},\mathbf{v}_{2},\ldots,\mathbf{v}_{n}$是线性无关的，当且仅当其张成$V$。
  \item 证明：向量空间$V$中的线性无关向量组$\mathbf{v}_{1},\mathbf{v}_{2},\ldots,\mathbf{v}_{n}$是基，当且仅当$n=\dim V$。
  \item （旧题重温：现在这道题应该很简单了）是否可能出现这样的情况：向量$\mathbf{v}_{1},\mathbf{v}_{2},\mathbf{v}_{3}$是线性相关的，但向量$\mathbf{w}_1=\mathbf{v}_1+\mathbf{v}_2,\mathbf{w}_2=\mathbf{v}_2+\mathbf{v}_3$与$\mathbf{w}_3=\mathbf{v}_3+\mathbf{v}_1$是线性{\bf 无关}的？{\bf 提示}：子空间$\operatorname{span}(\mathbf{v}_{1},\mathbf{v}_{2},\mathbf{v}_{3})$的维数是多少？
  \item 设向量组$\mathbf{u},\mathbf{v},\mathbf{w}$是$V$的基。证明：$\mathbf{u+v+w},\mathbf{v+w},\mathbf{w}$也是$V$的基。
  \item 考虑向量空间$\mathbb{R}^{5}$中的向量$\mathbf{v}_{1}=(2,-1,1,5,-3)^{T},\mathbf{v}_{2}=(3,-2,0,0,0)^{T},$ $\mathbf{v}_3=(1,1,50,-921,0)^T$。
  \begin{enumerate}[label={\alph*)}]
    \item 证明：这些向量是线性无关的；
    \item 将这个向量补足为基。
    
    如果你先解b)，那么你可以不作任何计算就完成整道题。
  \end{enumerate}
\end{problemset}

\section{线性方程组的通解}\label{sec:general_solution}

在接下来这短短一节中，我们讨论线性方程组的通解（即解集）的结构。

称方程组$A\mathbf{x}=\mathbf{b}$是齐次的，若右端项$\mathbf{b}=\mathbf{0}$，即齐次线性方程组形式为$A\mathbf{x}=\mathbf{0}$。

每个方程组\[A\mathbf{x}=\mathbf{b}\]都可通过取$\mathbf{b}=\mathbf{0}$关联到一个齐次方程组。

\begin{theorem}[线性方程组的通解]\label{thm:general_solution}
  令向量$\mathbf{x}_1$满足方程$A\mathbf{x}=\mathbf{b}$，并令$H$为其所关联的齐次方程组\[A\mathbf{x}=\mathbf{0}\]的解集。那么集合\[\{\mathbf{x}=\mathbf{x}_1+\mathbf{x}_\mathrm{h}:\mathbf{x}_\mathrm{h}\in H\}\]是方程$A\mathbf{x}=\mathbf{b}$的解集。
\end{theorem}

该定理还可陈述为
\begin{center}
  \fbox{$A\mathbf{x}=\mathbf{b}$的通解} = \fbox{$A\mathbf{x}=\mathbf{b}$的特解} + \fbox{$A\mathbf{x}=\mathbf{0}$的通解}
\end{center}
\begin{proof}
  取定向量$\mathbf{x}_{1}$满足$A\mathbf{x}_{1}=\mathbf{b}$。令向量$\mathbf{x}_\mathrm{h}$满足$A\mathbf{x}_\mathrm{h}=\mathbf{0}$。那么对于$\mathbf{x}=\mathbf{x}_{1}+\mathbf{x}_{\mathrm{h}}$有\[A\mathbf{x}=A(\mathbf{x}_1+\mathbf{x}_\mathrm{h})=A\mathbf{x}_1+A\mathbf{x}_\mathrm{h}=\mathbf{b}+\mathbf{0}=\mathbf{b},\]
  因此任意形如\[\mathbf{x}=\mathbf{x}_1+\mathbf{x}_\mathrm{h},\quad\mathbf{x}_\mathrm{h}\in H\]的$\mathbf{x}$都是$A\mathbf{x}=\mathbf{b}$的解。
  
  现设$\mathbf{x}$满足$A\mathbf{x}=\mathbf{b}$。那么对于$\mathbf{x}_\mathrm{h}:=\mathbf{x}-\mathbf{x}_1$有\[A\mathbf{x}_\mathrm{h}=A(\mathbf{x}-\mathbf{x}_1)=A\mathbf{x}-A\mathbf{x}_1=\mathbf{b}-\mathbf{b}=\mathbf{0},\]
  所以$\mathbf{x}_\mathrm{h}\in H$。因此$A\mathbf{x}=\mathbf{b}$的\textbf{任意}解$\mathbf{x}$都能表示为$\mathbf{x}=\mathbf{x}_1+\mathbf{x}_\mathrm{h}$，其中$\mathbf{x}_\mathrm{h}\in H$。
\end{proof}

这个定理的强大之处在于其普遍性。它适用于所有的线性方程，而且我们不需要假设向量空间是有限维的。你还会在微分方程、积分方程、偏微分方程等各种场合与这个定理重逢。除了表明解集的结构，这个定理还使我们得以将对唯一性的分析与对存在性的研究分开。也就是说，要研究唯一性，只需要分析齐次方程$A\mathbf{x}=\mathbf{0}$解的唯一性（这个方程总有解）。

就本书内容而言，这个定理有个直接的应用：我们可借以检验方程$A\mathbf{x}=\mathbf{b}$的解。
例如，考虑方程组
\[\begin{pmatrix}2&3&1&4&-9\\1&1&1&1&-3\\1&1&1&2&-5\\2&2&2&3&-8\end{pmatrix}\mathbf{x}=\begin{pmatrix}17\\6\\8\\14\end{pmatrix}.\]
作行变换，可求出该方程组的解
\begin{equation}\label{example_solution}
  \mathbf{x}=\begin{pmatrix}3\\1\\0\\2\\0\end{pmatrix}+x_3\begin{pmatrix}-2\\1\\1\\0\\0\end{pmatrix}+x_5\begin{pmatrix}2\\-1\\0\\2\\1\end{pmatrix},\quad x_3,x_5\in\mathbb{F}.
\end{equation}
此处的参数$x_3$和$x_5$可用任意其他字母替代，比如说$t$和$s$。用$x_3$和$x_5$来表示，只是为了提醒我们，这些参数来自相应的自由变量。

现在假设这个解已给出，并要检验其是否正确。当然，我们可以重做一次行变换，但这太费时了。并且，如果解不是通过标准的方法获得的，那么它可能与从行化简得到的解看起来不同。例如，
\begin{equation}\label{example_solution_2}
  \mathbf{x}=\begin{pmatrix}3\\1\\0\\2\\0\end{pmatrix}+s\begin{pmatrix}-2\\1\\1\\0\\0\end{pmatrix}+t\begin{pmatrix}0\\0\\1\\2\\1\end{pmatrix},\quad s,t\in\mathbb{F}
\end{equation}
与 \eqref{example_solution} 表示同样的解集（你能看出为什么吗？），这里我们只是将 \eqref{example_solution} 中的最后一个向量替换成了它和第二个向量之和。因此，虽然这个式子和我们根据行化简得到的解不同，但它是正确的。

检验 \eqref{example_solution} 和 \eqref{example_solution_2} 给出了正确解的最简单方法，就是检验第一个向量$(3,1,0,2,0)^T$满足方程$A\mathbf{x} = \mathbf{b}$，并且其余两个向量（前面有参数的那两个向量）满足原方程组所关联的齐次方程组$A\mathbf{x} = \mathbf{0}$。

如果这些检验都通过了，那么我们就可确信由 \eqref{example_solution} 或 \eqref{example_solution_2} 所确定的任意向量$\mathbf{x}$都是原方程的解。

注意，这种检验解的方法并不能保证 \eqref{example_solution} 或 \eqref{example_solution_2} 给出所有的解。例如，如果我们漏掉了系数为$x_3$的那一项，上面的检验仍会全部通过。

所以，如何保证在 \eqref{example_solution} 中，我们没有漏掉或是多出任何一个自由变量？

这就要重新数一下主元的数目。在本例中，作行变换可知主元数目为$3$。因此，的确应有$2$个自由变量，这样 \eqref{example_solution} 中就没有任何遗漏。

如果要{\bf 证明}解的正确性，我们就需要新的概念——矩阵的基本子空间和秩。还要指出，在本例中，没必要作完整的行变换来检验的确只有$2$个自由变量，\eqref{example_solution} 和 \eqref{example_solution_2} 也都是正确的通解表达式。

\begin{problemset}
  \item 判断正误：
  \begin{enumerate}[label={\alph*)}]
    \item 任意线性方程组都至少有一个解；
    \item 任意线性方程组都至多有一个解；
    \item 任意齐次线性方程组都至少有一个解；
    \item 任意具有$n$个方程$n$个未知数的线性方程组都至少有一个解；
    \item 任意具有$n$个方程$n$个未知数的线性方程组都至多有一个解；
    \item 如果对应于给定线性方程组的齐次方程组有解，那么原方程组有解；
    \item 如果具有$n$个方程$n$个未知数的齐次线性方程组的系数矩阵是可逆的，那么该方程组没有非零解；
    \item 任意具有$m$个方程$n$个未知数的线性方程组的解集都是$\mathbb{R}^n$的子空间；
    \item 任意具有$m$个方程$n$个未知数的齐次线性方程组的解集都是$\mathbb{R}^n$的子空间。
  \end{enumerate}
  \item 求出通解如下的$2\times 3$线性方程组（具有$2$个方程、$3$个未知数）。
  \[\begin{pmatrix}1\\1\\0\end{pmatrix}+s\begin{pmatrix}1\\2\\1\end{pmatrix},\quad s\in\mathbb{R}.\]
\end{problemset}

\section{矩阵的基本子空间、秩}

我们在第 \ref{chap:BasicNotion} 章第 \ref{sec:subspace} 节讨论过，每个线性变换$A:V\to W$都与两个子空间相关联，一个是它的核（零空间）\[\operatorname{Ker}A=\operatorname{Null}A:=\{\mathbf{v}\in V:A\mathbf{v}=\mathbf{0}\}\subset V,\]
另一个是它的值域\[\operatorname{Ran}A=\{\mathbf{w}\in W:\mathbf{w}=A\mathbf{v},\mathbf{v}\in V\}\subset W.\]
换言之，核$\ker A$是齐次方程组$A\mathbf{x}=\mathbf{0}$的解集，值域$\rg A$是使得方程$A\mathbf{x}=\mathbf{b}$有解的所有右端项$\mathbf{b}\in W$所成的集合。

若$A$是$m\times n$矩阵，即从$\mathbb{F}^{n}$到$\mathbb{F}^{m}$的映射，那么由矩阵-向量乘法的“列乘以坐标”法则，任意向量$\mathbf{w}\in\operatorname{Ran}A$可表示为$A$的列的线性组合。这就能解释{\bf 列空间}（记号为$\col A$）这一名称的来历，它是$\rg A$的另一种表达。

对于矩阵$A$，在$\rg A$和$\ker A$以外，还常讨论转置矩阵$A^T$的核与值域。通常，$\rg A^T$常称为{\bf 行空间}，而$\ker A^T$常称为{\bf 左零空间}（但往往不用特别的记号来表示它们）。

称$\operatorname{Ran}A,\operatorname{Ker}A,\operatorname{Ran}A^{T},\operatorname{Ker}A^{T}$这四个子空间是矩阵$A$的{\bf 基本子空间}。本节我们将研究四个基本子空间维数上的重要关系。

我们将会用到下面定义——它是线性代数中基础性的定义之一。
\begin{definition*}
  对于线性变换（矩阵）$A$，其秩$\rank A$是$A$的值域的维数\[\rank A:=\dim \rg A.\]
\end{definition*}

\subsection{计算基本子空间和秩}

为计算矩阵的基本子空间和秩，需要作行化简。也即，令$A$为矩阵，$A_{\mathrm{e}}$为其阶梯形
\begin{enumerate}
  \item \textbf{原}矩阵$A$的主元列（即作行变换之后，所得阶梯形中具有主元的那些列）就是$\rg A$的（一个）基。
  \item 阶梯形$A_{\mathrm{e}}$的主元{\bf 行}给出了行空间的一个基。当然，也可以取原矩阵的转置并作行变换。但是，如果我们已经得到$A$的阶梯形（比如说在计算$\rg A$的过程中），那么我们就可马上得到$\rg A^T$。
  \item 为了求出零空间$\ker A$的基，需要求解齐次方程$A\mathbf{x}=\mathbf{0}$：细节参见下面示例。
\end{enumerate}

\begin{example*}
  考虑矩阵\[\begin{pmatrix}1&1&2&2&1\\2&2&1&1&1\\3&3&3&3&2\\1&1&-1&-1&0\end{pmatrix}.\]
  作行变换可得阶梯形\[\begin{pmatrix}\boxed{1}&1&2&2&1\\0&0&\boxed{-3}&-3&-1\\0&0&0&0&0\\0&0&0&0&0\end{pmatrix}\]
  （这里框出了主元）。因此，{\bf 原矩阵}的第1、3列，也即\[\begin{pmatrix}1\\2\\3\\1\end{pmatrix},\quad\begin{pmatrix}2\\1\\3\\-1\end{pmatrix}\]
  给出了$\rg A$的一个基。我们还可马上得到行空间$\rg A^T$的基：{\bf 阶梯形}的第一、二行，即向量\[\begin{pmatrix}1\\1\\2\\2\\1\end{pmatrix},\quad\begin{pmatrix}0\\0\\-3\\-3\\-1\end{pmatrix}\]
  （此处我们将向量写成列向量。究竟是将向量像此处一样写成列向量还是写成行向量，是个习惯问题。这里我们写成列向量的理由是，尽管我们称$\rg A$是{\bf 行空间}，但我们其实将它定义为$A^T$的列空间）

  为计算零空间$\ker A$的基，我们需要解方程$A\mathbf{x}=\mathbf{0}$。计算$A$的{\bf 简化}阶梯形，对于本例是\[\begin{pmatrix}\boxed{1}&1&0&0&1/3\\0&0&\boxed{1}&1&1/3\\0&0&0&0&0\\0&0&0&0&0\end{pmatrix}.\]
  \pagebreak

  \noindent 注意，在求解齐次方程$A\mathbf{x}=\mathbf{0}$时，并不需要写出完整的增广矩阵。的确，此时增广矩阵的最后一列全为$0$，在行变换中不会改变。因此，我们不需要真的写出它，而只需记得有这一列，则可从上述简化阶梯形中看出解为\[\left\{\begin{array}{l}x_1=-x_2-\frac{1}{3}x_5,\\x_2~\text{是自由变量},\\x_3=-x_4-\frac{1}{3}x_5\\x_4~\text{是自由变量},\\x_5~\text{是自由变量},\end{array}\right.\]
  或写成向量形式
  \begin{equation}\label{basis_of_ker_a}
    \mathbf{x}=\begin{pmatrix}-x_2-\frac{1}{3}x_5\\x_2\\-x_4-\frac{1}{3}x_5\\x_4\\x_5\end{pmatrix}=x_2\begin{pmatrix}-1\\1\\0\\0\\0\end{pmatrix}+x_4\begin{pmatrix}0\\0\\-1\\1\\0\end{pmatrix}+x_5\begin{pmatrix}-1/3\\0\\-1/3\\0\\1\end{pmatrix}
  \end{equation}
  每个自由变量之后的向量，也即向量\[\begin{pmatrix}-1\\1\\0\\0\\0\end{pmatrix},\quad\begin{pmatrix}0\\0\\-1\\1\\0\end{pmatrix},\quad\begin{pmatrix}-1/3\\0\\-1/3\\0\\1\end{pmatrix}\]
  形成$\ker A$的基。

  遗憾的是，求出$\ker A^T$的基并没有捷径——必须求解方程$A^T\mathbf{x}=\mathbf{0}$。已知$A$的阶梯形并无帮助。
\end{example*}

\subsection{基本子空间的基的计算方式的解释}

所以，为什么以上方法确实能给出基本子空间的基呢？

\subsubsection{零空间$\ker A$}

零空间$\ker A$的基的求法解释起来可能最容易：因为我们求解了方程$A\mathbf{x}=\mathbf{0}$，即求出它所有的解，那么$\ker A$的所有向量都是我们所得各向量的线性组合。因此，我们所得各向量形成了$\ker A$的张成组。要看出这个向量组是线性无关的，让我们将每个向量都乘以对应的自由变量并求和——见 \eqref{basis_of_ker_a}。那么对于任意自由变量$x_k$，求和结果的第$k$项恰为$x_k$（同样可见 \eqref{basis_of_ker_a}），因此该结果（线性组合）为$\mathbf{0}$只有一种可能，就是所有自由变量都为$0$。

\subsubsection{列空间$\rg A$}

下面我们解释列空间$\rg A$的基的求法为何奏效。首先，注意到$A$的\textbf{简化阶梯形}$A_{\mathrm{re}}$的主元列形成$\rg A_{\mathrm{re}}$的基（并不是原矩阵的列空间的基，而是其简化阶梯形的基！）。因为行变换仅仅是左乘可逆矩阵，所以它并不改变线性无关关系。因此，{\bf 原}矩阵$A$的主元列也是线性无关的。

现在我们说明$A$的主元列张成$A$的列空间。令$\mathbf{v}_{1},\mathbf{v}_{2},\ldots,\mathbf{v}_{r}$是$A$的主元列，并令$\mathbf{v}$是$A$的任意一列。我们要证明：$\mathbf{v}$可表示成主元列$\mathbf{v}_{1},\mathbf{v}_{2},\ldots,\mathbf{v}_{r}$的线性组合\[\mathbf{v}=\alpha_1\mathbf{v}_1+\alpha_2\mathbf{v}_2+\cdots+\alpha_r\mathbf{v}_r.\]
简化阶梯形$A_{\mathrm{re}}$是由$A$左乘初等矩阵得到的：\[A_{\mathrm{re}}=EA\]
其中$E$是一系列初等矩阵的积，因此$E$是可逆矩阵。向量$E\mathbf{v}_{1},E\mathbf{v}_{2},\ldots,$ $E\mathbf{v}_{r}$是$A_{\mathrm{re}}$的主元列，而$A$的其中一列$\mathbf{v}$被变换为$A_{\mathrm{re}}$的一列$E\mathbf{v}$。因为$A_{\mathrm{re}}$的主元列形成$\rg A_{\mathrm{re}}$的基，所以向量$E\mathbf{v}$可表示为它们的线性组合\[E\mathbf{v}=\alpha_1E\mathbf{v}_1+\alpha_2E\mathbf{v}_2+\ldots+\alpha_rE\mathbf{v}_r.\]
将该等式两端同时左乘$E^{-1}$，可得如下表达式\[\mathbf{v}=\alpha_1\mathbf{v}_1+\alpha_2\mathbf{v}_2+\cdots+\alpha_r\mathbf{v}_r,\]
因此$A$的主元列的确张成$\rg A$。

\subsubsection{行空间$\rg A^T$}

容易看出，$A$的\textbf{阶梯形}$A_{\mathrm{e}}$的主元行是线性无关的。的确如此：令$\mathbf{w}_1,\mathbf{w}_2,\ldots,\mathbf{w}_r$是$A_{\mathrm{e}}$的转置的主元行（因为我们总将这些向量写成列向量）。假设\[\alpha_1\mathbf{w}_1+\alpha_2\mathbf{w}_2+\ldots+\alpha_r\mathbf{w}_r=\mathbf{0}.\]
考虑$\mathbf{w}_1$的第一个非零项。因为其他所有向量$\mathbf{w}_2,\mathbf{w}_3,\ldots,\mathbf{w}_r$的对应项都等于$0$（由阶梯形的定义），我们可得$\alpha_1=0$。因此可以直接忽略求和式的第一项。

现在考虑$\mathbf{w}_2$的第一个非零项。向量$\mathbf{w}_3,\ldots,\mathbf{w}_r$的对应项都为$0$，因此$\alpha_2=0$。重复此过程，我们可得$\alpha_{k}=0\ \forall k=1,2,\ldots,r$。

要看出$\mathbf{w}_{1},\mathbf{w}_{2},\ldots,\mathbf{w}_{r}$张成行空间，请注意{\bf 行变换并不改变行空间}。直接分析行变换过程就能看出这一点，不过此处我们给出一种更为正式的证法。

对于变换$A$和集合$X$，我们记$A(X)$是所有可被表示成$y=A(x),$ $x\in X$的元素$y$所成的集合\[A(X):=\{y=A(x):x\in X\}.\]

若$A$是$m\times n$矩阵，$A_{\mathrm{e}}$是其阶梯形，$A_{\mathrm{e}}$是由$A$左乘初等矩阵得到的：\[A_{\mathrm{e}}=EA\]
其中$E$是$m\times m$可逆矩阵（初等矩阵的积）。那么\[\operatorname{Ran}A_{\mathrm{e}}^T=\operatorname{Ran}(A^TE^T)=A^T(\operatorname{Ran}E^T)=A^T(\mathbb{R}^m)=\operatorname{Ran}A^T,\]
因此，的确有$\rg A^T = \rg A_{\mathrm{e}}^T$。

\subsection{秩定理、基本子空间的维数}

在许多实际应用中，都需要求出矩阵列空间或零空间的基。例如，如上所示，求解齐次方程组$A\mathbf{x}=0$就是求出零空间$\ker A$的基。求出列空间的基也就意味着从张成组中找出一个基（通过移除不必需的列向量）。

然而，上述基本子空间的基的计算方法的最重要应用，就是得出这些子空间的维数关系。

\begin{theorem}[秩定理]
  对于矩阵$A$，\[\rank A = \rank A^T.\]
\end{theorem}
该定理常表述为：
\begin{center}
  \fbox{矩阵的{\bf 列秩}等于{\bf 行秩}。}
\end{center}
该定理的证明很简单，因为$\rg A$和$\rg A^T$的维数都等于$A$的阶梯形的主元的个数。

下面定理给出了基本子空间的维数之间的重要关系，它也常被称为秩定理。

\begin{theorem}\label{thm:fundamental_thm_of_la}
  令$A$是$m\times n$矩阵，即从$\mathbb{F}^n$到$\mathbb{F}^m$的线性变换。那么
  \begin{enumerate}
    \item $\dim \ker A + \rank A = n$（$A$的定义域的维数）；
    \item $\dim \ker A^T + \rank A = m$（$A$的目标空间的维数）。
  \end{enumerate}
\end{theorem}
\begin{proof}
  利用上述基本子空间的基的求法，本证明非常简单。第一个结论就是，自由变量的数目（$\dim\ker A$）加上基本变量的数目（主元的数目，即$\rank A$）等于列数（$n$）。

  而第二个结论，只需注意到$\rank A = \rank A^T$，就可知道仅仅是将第一个结论应用于$A^T$的结果。
\end{proof}

为演示上述定理的应用，考虑第 \ref{sec:general_solution} 节的例子。彼时我们考虑了线性方程组\[\begin{pmatrix}2&3&1&4&-9\\1&1&1&1&-3\\1&1&1&2&-5\\2&2&2&3&-8\end{pmatrix}\mathbf{x}=\begin{pmatrix}17\\6\\8\\14\end{pmatrix}.\]
并求出了其通解表达式
\begin{equation*}
  \mathbf{x}=\begin{pmatrix}3\\1\\0\\2\\0\end{pmatrix}+x_3\begin{pmatrix}-2\\1\\1\\0\\0\end{pmatrix}+x_5\begin{pmatrix}2\\-1\\0\\2\\1\end{pmatrix},\quad x_3,x_5\in\mathbb{F}.
\end{equation*}
或
\begin{equation*}
  \mathbf{x}=\begin{pmatrix}3\\1\\0\\2\\0\end{pmatrix}+s\begin{pmatrix}-2\\1\\1\\0\\0\end{pmatrix}+t\begin{pmatrix}0\\0\\1\\2\\1\end{pmatrix},\quad s,t\in\mathbb{F}
\end{equation*}
在第 \ref{sec:general_solution} 节中我们检验了这两个式子给出的向量$\mathbf{x}$都确实是原方程的解。但是，我们如何保证这两个式子给出了{\bf 所有}解呢？

首先，我们知道这两个式子中，最后$2$个向量（前面乘了参数的向量）属于$\ker A$。容易看出，两式中的这两个向量都是线性无关的（两个向量线性相关，当且仅当其中一个向量是另一个的若干倍）。

现在计算维数：交换前两行并作一轮行变换：
\begin{align*}
  \begin{pNiceArray}{rrrrr}[margin,first-col]
    &1&1&1&1&-3\\
    -2R_1&2&3&1&4&-9\\
    -R_1&1&1&1&2&-5\\
    -2R_1&2&2&2&3&-8
  \end{pNiceArray} &\sim 
  \begin{pNiceArray}{rrrrr}[margin]
     1&1&1&1&-3\\0&1&-1&2&-3\\0&0&0&1&-2\\0&0&0&1&-2
  \end{pNiceArray}
\end{align*}
从中已能看出有三个主元，因此$\rank A\ge 3$。（其实我们已经能看出秩为$3$，但是此处只需要推到这个不等式就够了）。由\autoref{thm:fundamental_thm_of_la}，$\rank A+\dim\ker A = 5$，于是$\dim\ker A\le 2$，因此$\ker A$中不会有超过$2$个向量线性无关。因此，上面表达式中的最后两个向量都形成$\ker A$的基，则这两个表达式都给出了原方程的所有解。

秩定理的一个重要推论，就是下面这个将线性方程组解的存在性和唯一性关联起来的定理。

\begin{theorem}
  令$A$是$m\times n$矩阵。那么方程\[A\mathbf{x}=\mathbf{b}\]
  对于任一$\mathbf{b}\in\mathbb{R}^m$都有解，当且仅当对偶方程\[A^T\mathbf{x}=\mathbf{0}\]都具有唯一解（仅有平凡解）。（注意，第二个方程中是$A^T$，不是$A$）
\end{theorem}
\begin{proof}
  由\autoref{thm:fundamental_thm_of_la} 可立即得证（只需计算维数）。我们将证明的细节留给读者作为习题。
\end{proof}

第二个秩定理（\autoref{thm:fundamental_thm_of_la}）有一个非常优美的几何解释：该定理的第一条是说，如果变换$A:\mathbb{F}^{n}\to\mathbb{F}^{m}$具有平凡的核（$\operatorname{Ker}A=\{\mathbf{0}\}$），那么定义域$\mathbb{F}^{n}$和值域$\rg A$的维数相同。如果核不是平凡的，那么该变换会“消掉”$\dim\ker A$维，从而$\dim\rg A = n-\dim\ker A$。

\section{线性变换在任意基下的表示、坐标变更公式}\label{sec:ChangeOfCoordinate}

% 如下是算法~\ref{alg:euclid}.
% \begin{algorithm}[H]
%   \small
%   \caption{~Euclid's algorithm}\label{alg:euclid}
%   \begin{algorithmic}[1]
%     \Procedure{Euclid}{$a,b$}\Comment{The g.c.d. of a and b}
%     \State $r\gets a\bmod b$
%     \While{$r\not=0$}\Comment{We have the answer if r is 0}
%     \State $a\gets b$
%     \State $b\gets r$
%     \State $r\gets a\bmod b$
%     \EndWhile\label{euclidendwhile}
%     \State \Return $b$\Comment{The gcd is b}
%     \EndProcedure
%   \end{algorithmic}
% \end{algorithm}


% 如下是算法~\ref{alg:foo}, 算法宽度可以通过 minipage 宏包调节.
% \begin{center}
%   \vspace{-2ex}
%   \begin{minipage}{.9\linewidth}
%     \begin{algorithm}[H]
%       \caption{~算法的名字}\label{alg:foo}
%       \begin{algorithmic}[1]
%         \Require input parameters A, B, C
%         \Ensure output result
%         \State some description 算法介绍
%         \For{condition}
%         \State ...
%         \If{condition}
%         \State ...
%         \Else
%         \State ...
%         \EndIf
%         \EndFor
%         \While{condition}
%         \State ...
%         \EndWhile
%         \State \Return result
%       \end{algorithmic}
%     \end{algorithm}
%   \end{minipage}
% \end{center}


% \section{\texorpdfstring{{\boldmath$Z=X \cup Y$}}{Z = X union Y} 的情况}

% 这是一个小节标题中出现数学符号的情况.

